\documentclass[11pt,a4paper]{article}
\usepackage[T1]{fontenc}
\usepackage{isabelle,isabellesym}
\usepackage{amsmath,amssymb}
\usepackage{pdfsetup}

\urlstyle{rm}
\isabellestyle{it}

\begin{document}
\emergencystretch=3em

\title{Regulatory Action Composition}
\author{Jinwook Kim (Jay)}
\maketitle

\begin{abstract}
We mechanize sequential regulatory-action composition over a concrete
five-state, seven-label reference machine in Isabelle/HOL.  The theory
separates applied transitions, legal rejection, and operational failure;
classifies all 21 unordered pairs of distinct transition labels into 12
commuting and 9 noncommuting pairs with checked witnesses; proves terminal
absorption and restrictive-state provenance; specifies typed legal effects
and state-sensitive transfer gates; records trace and frame properties for
commands and completed atomic synchronization steps; and enumerates exactly
60 distinct reachable state-transformation vectors.  These are model-level
results for the declared finite machine.  They do not establish legal truth,
authorization correctness, concurrent execution, distributed rollback,
network liveness, or model-to-code refinement.
\end{abstract}

\tableofcontents

\section{Overview}

Regulatory commands do not in general commute.  Applying two otherwise valid
labels in opposite orders may produce different final states because the
first transition changes whether the second transition is defined.  This
artifact makes that order sensitivity explicit and executable at the model
level, while preserving a distinction between an undefined legal transition
and an operational failure outside the state machine.

The session is a standalone artifact built on the published
\texttt{Cross\_Domain\_State\_Preservation} session.  It imports only the
concrete regulatory instance from that session; it does not import or extend
the authenticated-data-structure functor layer.  The shared base supplies the
five states, seven transition labels, transition function, and atomic global
state abstraction~\cite{cdsp}.  The regulatory interpretation is motivated by
the protocol model described in the related RCP paper~\cite{rcp}.

\subsection{Checked result groups}

\begin{description}
\item[Outcome semantics.] Applied, rejected, and operational-failure outcomes
  remain distinct, with rejected and failed commands leaving the modeled
  state unchanged.
\item[Pair classification.] The complete set of 21 unordered pairs of
  distinct labels is partitioned into 12 commuting and 9 noncommuting pairs;
  every noncommuting pair has a concrete state witness.
\item[Terminality and provenance.] Confiscation is absorbing for arbitrary
  action words, while the reachable restrictive states carry checked
  provenance constraints.
\item[Legal-effect and transfer layer.] Six typed legal-effect descriptors
  cover state transitions and the separate enforcement-transfer path.
  Ordinary transfers use both the regulatory-state gate and a current-policy
  condition.
\item[Trace and synchronization.] Command queues preserve order and input
  identity, record execution outcome, leave unrelated assets unchanged, and
  preserve validity after each completed atomic synchronization step.
\item[Finite normal forms.] Exactly 60 distinct state-transformation vectors
  are reachable; every action word evaluates to one of them and the set is
  closed under every transition label.
\end{description}

\subsection{Assurance boundary}

The artifact is a finite reference semantics, not a legal opinion or an
implementation proof.  In particular, the current-policy input for restricted
transfers is abstract.  Asset lookup, authorization, replay protection,
concurrent interleavings, partial propagation, timeout, rollback, bytecode
behavior, and deployment security remain external or refinement obligations.
The atomic queue theorems speak only about completed model steps.

\input{session}

\section{Conclusion}

The checked algebra shows where regulatory actions are order-insensitive and
where order is semantically load-bearing.  It also supplies a bounded,
auditable normal-form universe for regression tests and later implementation
mapping, without extending the claims beyond the concrete reference machine.

\bibliographystyle{plain}
\bibliography{root}

\end{document}
